\section{Composition}

{
  \color{blue}
  \section{Todo}

  % failure/counter example
  % 3 failure test cases in composition
  % 21 failure test cases pairwise
  % add back fafilure senario in XC

  % principal of how to model apps (when to use symbol, when to use
  % instantiation)
  When modeling systems for consistency analysis, choosing the right
  approach to abstract away underlying system behavior is crucial.
  We propose the following principles to guide the modeling process:

  % using default symbols
  \textbf{Default symbols} are appropriate when system has straightforward
  composition patterns with minimal cross-service dependencies.
  Services should have well-defined consistency guarantees that don't
  require fine-grained differentiation.
  The primary concern is with verifying that the overall system
  composition satisfies basic consistency requirements.

  As shown in the online shopping mall example, default symbols work well when
  modeling a system where each component (shop, transaction log,
  arbitrator) has clear, isolated responsibilities. The focus is on
  ensuring that the composition of these components maintains the
  desired consistency guarantees without needing to track specific data
  objects or operation instances.

  % operation instantiation
  \textbf{Operation instantiations} should be used
  when the system's consistency depends on the ordering of specific API
  calls or RPC functions.
  This is most useful when there are ordering constraints on API/RPC
  calls or dependencies between actions (as in flow of data) across services
  The focus is on the behavior of the system rather than the underlying
  data storage.

  The media example illustrates this approach by modeling each API
  action as a distinct operation instance. This allows specification of
  ordering constraints between operations,
  such as requiring login before accessing content or checking metadata
  before serving videos.

  % custom symbols

  Our current framework abstracts system-wide behavior from a single
  node's perspective, applying constraints to operations primarily within the
  scope of a single history and its corresponding abstract
  executions. While this
  focuses the analysis and simplifies reasoning on local consistency properties
  from the observing node's perspective, it also means that global behaviors
  across multiple nodes need to be modeled separately on different
  set of symbols
  to avoid conflicts in quantified variables. Further, to reason about cross
  system behaviors, ``glue'' constraints need to be added to represent data flow
  between different systems.

  \textbf{Custom symbols} become necessary when different storage
  systems maintain separate state spaces that need to
  be distinguished between each other, such as cross service ordering
  constraints exist between specific data objects.
  Or we need to track how consistency properties apply to particular data
  elements across system boundaries. This means that each consistency model
  object instantiation will be different even though the same model is used.
  This is required as the system's internal behavior are captured using one set
  of symbolic representations, while the external behavior is captured using
  another. When interactions are required between internal and external, a shim
  layer is needed to translate between the two representations (i.e.
  symbolic equivalence in cross-causal).

  The cross-causal (XC) example demonstrates this approach, where each
  storage system has its own state and the consistency semantics must
  be explicitly differentiated using custom symbols.
}
